% !TeX program = xelatex

% compile with XeLaTeX
% this template was created by salim bou 
\documentclass[dvipsnames,mathserif, aspectratio=169]{beamer}
\usepackage{setspace, float}
\setstretch{1.2}
\usepackage{tikz}
\usepackage{subcaption, tasks}

\usepackage{amsmath, amssymb, amsthm, nicematrix, pst-node, nccmath, cancel}

\AtBeginDocument{
	\allowdisplaybreaks
	\renewcommand{\jot}{10pt}
	\abovedisplayskip=5pt
	\belowdisplayskip=5pt
	\renewcommand{\theequation}{\arabic{equation}}
}

\usepackage{polyglossia, quran}
\setdefaultlanguage[numerals=maghrib]{arabic} % locale=mashriq, libya, algeria, tunisia, morocco, or mauritania  for names of months in \date 
\setotherlanguage{english}
\newfontfamily\arabicfont[Script=Arabic]{Amiri}
\newfontfamily\arabicfontsf[Script=Arabic]{Amiri}
\newfontfamily{\timesfont}{Times New Roman}

\newcommand{\ar}{\textarabic}
\newcommand{\en}{\textenglish}

\usepackage[T1]{fontenc}
\usepackage{times}

\usepackage[lite, zswash]{mtpro2}

\DeclareMathSymbol{0}{\mathalpha}{operators}{`0}
\DeclareMathSymbol{1}{\mathalpha}{operators}{`1}
\DeclareMathSymbol{2}{\mathalpha}{operators}{`2}
\DeclareMathSymbol{3}{\mathalpha}{operators}{`3}
\DeclareMathSymbol{4}{\mathalpha}{operators}{`4}
\DeclareMathSymbol{5}{\mathalpha}{operators}{`5}
\DeclareMathSymbol{6}{\mathalpha}{operators}{`6}
\DeclareMathSymbol{7}{\mathalpha}{operators}{`7}
\DeclareMathSymbol{8}{\mathalpha}{operators}{`8}
\DeclareMathSymbol{9}{\mathalpha}{operators}{`9}

\newcommand{\N}{\mathbb{N}}
\newcommand{\Z}{\mathbb{Z}}
\newcommand{\Q}{\mathbb{Q}}
\newcommand{\R}{\mathbb{R}}
\newcommand{\C}{\mathbb{C}}
\newcommand{\LL}{\mathcal{L}}

\usetheme{Warsaw}
\setbeamertemplate{bibliography item}{\insertbiblabel}
%\usecolortheme{crane}

% for RTL liste
\makeatletter
\newcommand{\RTListe}{\raggedleft\rightskip\leftm}
\newcommand{\leftm}{\@totalleftmargin}
\makeatother



% RTL frame title
\setbeamertemplate{frametitle}
{\vspace*{-1mm}
	\nointerlineskip
	\begin{beamercolorbox}[sep=0.3cm,ht=2.2em,wd=\paperwidth]{frametitle}
		\vbox{}\vskip-2ex%
		\strut\hskip1ex\insertframetitle\strut
		\vskip-0.8ex%
	\end{beamercolorbox}
}


% align subsection in toc
\makeatletter
\setbeamertemplate{subsection in toc}
{\leavevmode\rightskip=5ex%
	\llap{\raise0.1ex\beamer@usesphere{subsection number projected}{bigsphere}\kern1ex}%
	\inserttocsubsection\par%
}
\makeatother

% RTL triangle for itemize
\setbeamertemplate{itemize item}{\scriptsize\raise1.25pt\hbox{\donotcoloroutermaths$\blacktriangleleft$}} 

%\setbeamertemplate{itemize item}{\rule{4pt}{4pt}}

\defbeamertemplate{enumerate item}{square2}
{\LR{
		%
		\hbox{%
			\usebeamerfont*{item projected}%
			\usebeamercolor[bg]{item projected}%
			\vrule width2.25ex height1.85ex depth.4ex%
			\hskip-2.25ex%
			\hbox to2.25ex{%
				\hfil%
				{\color{fg}\insertenumlabel}%
				\hfil}%
		}%
}}

\setbeamertemplate{enumerate item}[square2]

\setbeamertemplate{navigation symbols}{}

\let\definition\relax
\let\enddefinition\relax

\newtheorem{definition}{تعريف}

\let\example\relax
\let\endexample\relax
\newtheorem{example}{مثال}





\begin{document}
	\author{\textbf{الطالبة : عذراء عذاب عدنان}}
	\title{\textbf{تأثير الزمر على المجموعات}}
	\date{\textbf{اشراف\\
			 م. صفاء عبد الشهيد}}
	
		
		\begin{frame}
			\maketitle
		\end{frame}
		
		\begin{frame}{الخلاصة}
			\begin{itemize}
				\item \textbf{الفصل الأول:} في هذا الفصل تتطرقنا الى اهم المفاهيم في نظرية الزمر منها مفهوم الزمرة الزمرة الجزئية والتشاكل الزمري ، ....\\
				\item \textbf{الفصل الثاني:} في هذا الفصل عرّفنا مفهوم تأثير الزمر على المجموعات وأخذنا بعض الأمثلة وكذلك تعرّفنا على بعض النتائج المهمة\\
				\item \textbf{الفصل الثالث:} في هذا الفصل وسّعنا مفهوم نظرية كيلي في الزمر وطبقّنا النظرية بمثال.
			\end{itemize}
		\end{frame}
		
		\begin{frame}
			\begin{center}
				\Huge
				\textbf{الفصل الأول}\\
				\textbf{مفاهيم أساسية}
			\end{center}
		\end{frame}
		
		\begin{frame}
			\begin{definition}[الزمرة]
				يكون النظام الرياضي $(G, *)$ حيث $G$ مجموعة غير خالية و $*$ عملية ثنائية على $G$ زمرة (group) إذا تحقق ما يلي:
				\begin{enumerate}
					\item المجموعة $G$ مغلقة تحت العملية $*$، أي أن:
					\[
					\forall a, b \in G \Rightarrow a * b \in G
					\]
					
					\item العملية $*$ تحقق خاصية التجميع، أي أن:
					\[
					\forall a, b, c \in G \Rightarrow a * (b * c) = (a * b) * c
					\]
					
					\item يوجد عنصر محايد $e \in G$ بحيث:
					\[
					a * e = e * a = a \quad \forall a \in G
					\]
					
					\item لكل عنصر $a \in G$ يوجد عنصر معاكس $b \in G$ بحيث:
					\[
					a * b = b * a = e
					\]	
				\end{enumerate}
			\end{definition}
		\end{frame}
		
		\begin{frame}
			\begin{definition}[{الزمرة الجزئية}]
				لتكن $(G, *)$ زمرة، ولتكن $H$ مجموعة جزئية غير خالية من $G$. تُسمى $H$ زمرةً جزئية من $G$ إذا تحققت الشروط التالية:
				\begin{enumerate}
					\item المجموعة $H$ مغلقة تحت العملية $*$، أي أن:
					\[
					\forall a, b \in H \Rightarrow a * b \in H
					\]
					
					\item لكل عنصر $a \in H$ يوجد معكوسه في $H$، أي أن:
					\[
					\forall a \in H \Rightarrow a^{-1} \in H
					\]
				\end{enumerate}
			\end{definition}
		\end{frame}
		
		\begin{frame}
			\begin{definition}[{الزمرة الجئية الناظمية}]
				لتكن $(G, *)$ زمرة و $H$ زمرة جزئية من $G$. تُسمى $H$ زمرةً ناظمية في $G$ إذا تحقق الشرط:
				\[
				a * H = H * a \quad \forall a \in G
				\]
				حيث:
				\[
				a * H = \{a * h \mid h \in H\}, \quad H * a = \{h * a \mid h \in H\}.
				\]
				ويرمز لها بـ $H \trianglelefteq G$.
			\end{definition}
		\end{frame}
		
		\begin{frame}
			\begin{center}
				\Huge
				\textbf{الفصل الثاني}\\
				\textbf{تأثير الزمر على المجموعات}
			\end{center}
		\end{frame}
		
		\begin{frame}
			\begin{definition}[{تأثير الزمر على المجموعات}]
				لتكن $G$ زمرة ولتكن $A$ مجموعة غير خالية ، يعرف تأثير الزمرة $G$ على المجموعة $A$
				\linebreak(\en{action of group on a set})
				بأنه تطبيق 
				$* : G \times A \to A$ ( في العادة نكتب $g*a$ بدلاً من $*(g, a)$ ) الذي يحقق ما يلي:
				\begin{enumerate}
					\item $e * a = a$ لكل $a\in A$.
					\item $(g_1g_2)*a = g_1 * (g_2 * a)$ لكل $a\in A$ ولكل $g_1,g_2 \in G$.
				\end{enumerate} 
			\end{definition}
			
			\begin{example}
				اذا كان $*:G\times A \to A$ معرفاً بالقاعدة 
				$g*a = a$ لكل $g\in G$ وكل $a\in A$ فيمكن اعتبار ان $G$ تؤثر على $A$ لأن:
				\begin{enumerate}
					\item $e*a=a$ لكل $a\in A$.
					\item $(g_1g_2)*a=a=g_2*a=g_1*(g_2*a)$. لكل $g_1g_2\in G$ ولكل $a\in A$.
				\end{enumerate}	
				يسمى هذا التأثير بالتأثير التافه.
			\end{example}
		\end{frame}
		
		\begin{frame}
			\begin{example}
				اذا كانت $G$ هي زمرة التبديلات على المجموعة $A$ فأن 
				$*:G\times A \to A$ بالقاعدة
				$\sigma * a = \sigma(a)$ لكل $\sigma\in G$ ولكل $a\in A$ تأثيراً للزمرة $G$ على $A$ لأن
				\begin{enumerate}
					\item $e * a = e(a) = a$ لكل $a \in A$.
					\item $(\sigma_1 \circ \sigma_2)*a = (\sigma_1\circ\sigma_2)(a) = \sigma_1(\sigma_2(a)) = \sigma_1*(\sigma_2(a)) = \sigma_1 * (\sigma_2 * a)$. لكل 
					$\sigma_1, \sigma_2 \in G$ ولكل $a\in A$.
				\end{enumerate}
			\end{example}
		\end{frame}
		
		\begin{frame}
			\begin{theorem}\label{thm:actionrelation}
				لنفرض ان الزمرة $G$ تؤثر على المجموعة $A$ ، اذا كانت $\approx$ علاقة معرفة على $A$ كالتالي:\linebreak
				$a\approx b \iff \exists g\in G(ga= b)$ لكل $a, b\in A$ فأن $\approx$ علاقة تكافؤ على $A$.
			\end{theorem}
			\begin{definition}[{المدارات}]
				لتكن الزمرة $G$ تؤثر على المجموعة غير الخالية $A$.
				\begin{enumerate}
					\item تسمى صفوف تكافؤ العلاقة $\approx$ في المبرهنة \ref{thm:actionrelation} ، مدارات (\textbf{orbits}) $G$ على $A$.
					\linebreak لاحظ أن صف التكافؤ الذي يحتوي $a\in A$ هو:
					\[
					[a] = \{ b\in A : \exists g\in G \,\,( ga = b)\} = \{ga : g\in G\}
					\]
					\item يكون تأثير $G$ على $A$ متعدياً (transitive) اذا كان هناك مدار واحد فقط اي انه لكل $a, b\in A$ يوجد $g\in G$ حيث 
					$ga=b$.
				\end{enumerate}
			\end{definition}
		\end{frame}
		
		\begin{frame}
			\begin{center}
				\Huge
				\textbf{الفصل الثالث}\\
				\textbf{تعميم مبرهنة كيلي}
			\end{center}
		\end{frame}
		
		\begin{frame}
			في هذا الفصل سوف نناقش تطبيق مهم لتأثير الزمرة على المجموعة وهو تعميم لمبرهنة كيلي وكذلك بعض النتائج على هذه المبرهنة منها مبرهنة الدليل.
			
			\begin{theorem}
				لتكن $G$ زمرة وليكن $a\in G$
				\begin{enumerate}[label=(\alph*)]
					\item التطبيق $\lambda_a : G \to G$ المعرف بالقاعدة $\lambda_a(x) = ax$ لكل $x\in G$ تبديلاً.
					
					\item $H = \{\lambda_a : a\in G\} \le S_G$.
					
					\item $G \simeq H$.  
				\end{enumerate}
			\end{theorem}
		\end{frame}
		
		\begin{frame}
			\begin{theorem}
				اذا اثرت الزمرة $G$ على المجموعة $A$ بالضرب من اليسار (اي ان
				$g*a=ga$) فأن:
				\begin{enumerate}[label=(\alph*)]
					\item لكل $g\in G$ يكون التطبيق $\tau_g : A \to A$ المعرف بالقاعدة $\tau_g(a)= ga$ تبديلاً.
					\item $\tau_{g_1g_2} = \tau_{g_1} \circ \tau_{g_2}$ لكل $g_1, g_2\in G$.
					\item يكون التطبيق $\psi : G \to S_A$ المعرف بالقاعدة $\psi(g) = \tau_g$.
				\end{enumerate}
			\end{theorem}
			
			\begin{theorem}[مبرهنة كيلي المعممة]
				\label{thm:general}
				\addcontentsline{toc}{section}{مبرهنة (3-3)}
				اذا كانت $H\le G$ وكانت $A = \{aH : a\in G\}$ فأنه يوجد تشاكل $\psi : G\to S_A$ يحقق
				$\ker\psi \subseteq H$
			\end{theorem}
			
			\begin{theorem}[مبرهنة الدليل]
				\addcontentsline{toc}{section}{مبرهنة (3-4)}
				اذا كانت $H$ زمرة جزئية من الزمرة المنتهية $G$ حيث
				$[G : H] =n$ ، واذا كان $|G|$ لا يقسم $n!$
				\linebreak فإن $G$ تحتوي على زمرة جزئية ناظمية غير تافهة.
			\end{theorem}
		\end{frame}
		
		\begin{frame}
			\begin{center}
				\Huge
				\textbf{شكراً لحسن استماعكم}
			\end{center}
		\end{frame}
\end{document}